%------------------------------------------------------------------------------%
\section{Proposição, Teoremas}

Uma \emph{Proposição} é uma oração afirmativa, cujo valor lógico é falso ou
verdadeiro, sem dar espaço para uma terceira alternativa. Considere, como
exemplo, as seguintes proposições:
\begin{enumerate}[label=\alph*)]
    \item Todo triângulo é retângulo;
    \item A soma de todos os ângulos de um triângulo é \ang{180};
    \item A soma de dois números pares é sempre um número par;
    \item Todo Brasileiro é carioca;
    \item Se $a>1$ então $a^2<a$.
\end{enumerate}
Note que as afirmações (a), (d) e (e) são falsas. Para verificar que, de
fato, que uma proposição é falsa, utilizamos o que chamamos de
\emph{contraexemplos}. \emph{Contraexemplos} são casos que contrariam uma
proposição, ou seja, um exemplo que refuta uma determinada afirmação universal.

Um triângulo cujos ângulos interno são \ang{120}, \ang{40} e
\ang{20} não é um triângulo retângulo, uma vez que não possui um ângulo
interno de \ang{90}. Logo esse contra-exemplo refuta a proposição (a)
mostrando que tal afirmação é falsa. Para verificar a veracidade da afirmação
(d) teríamos que consultar o registro de todos os brasileiros e verificar se
nasceu no Rio de Janeiro, mas isto é falso, pois o conhecido escritor Graciliano
Ramos é um brasileiro nascido em Alagoas. Do mesmo modo, se escolhermos $a=2$,
podemos ver que $a^2=4 > a$ e, portanto, esse contraexemplo contradiz a
afirmação feita em (e), mostrando que essa também é falsa.

Por outro lado, as proposições (b) e (c) são verdadeiras e para nos
convencermos precisaríamos demonstrar usando argumentos matemáticos. Adiante
falaremos mais um pouco sobre isso. Observe que a proposição (e) é do tipo
``Se $P$, então $Q$'', onde $P$ e $Q$ são sentenças. No item (e), por exemplo,
temos que
\begin{gather*}
  P\colon a>1 \\
  Q\colon a^2<a.
\end{gather*}
Isso significa que estamos assumindo que $P$ é verdade e usando este fato
devemos verificar se $Q$ é verdade ou não. Uma \emph{proposição condicional}
ou \emph{implicativa} é uma nova proposição formada a partir de duas
proposições P e Q, que é escrita na forma:
\[
  \text{Se }       P,
  \text{ então }   Q
  \text{ ou }      P
  \text{ implica } Q
\]
Denotamos a expressão
$P$ implica $Q$ por $P \Rightarrow Q$

No caso, a sentença P é chamada de \emph{hipótese} e a sentença $Q$ é
denominada de \emph{tese} e a validade da hipótese implica na veracidade da
tese.

A \emph{recíproca} de uma proposição condicional ou implicativa é a
proposição:
\[
  \text{Se }     Q,
  \text{ então } P
  \text{ ou }    Q \Rightarrow P
\]
A \emph{Contrapositiva} ou \emph{Contrarecíproca}
de uma afirmação condicional é a expressão :
\[
  \text{Se }     \sim Q,
  \text{ então } \sim P
  \text{ ou }    \sim Q
  \Rightarrow    \sim P
\]
Vale ressaltar que a contrapositiva tem o mesmo valor lógico da proposição a
ela associada, ou seja, dizer que $P \Rightarrow Q$ é o mesmo que dizer
$\sim Q \Rightarrow \sim P$.
Vamos chamar o modo em enunciamos uma proposição de
\emph{forma  positiva}.

Para que fique claro, considere a seguinte proposição:
\[
  \text{Se como laranja então gosto de fruta.}
\]
Na proposição enunciada acima temos que a hipótese é ``como laranja''
e a tese é ``gosto de frutas''.
Por outro lado a recíproca dessa proposição é
\[
  \text{Se gosto de frutas então como laranja.}
\]

Por fim, a contrapositiva é
\[
  \text{Se não gosto de frutas então não como laranja.}
\]

Mais a frente veremos que a forma contrapositiva de uma proposição poderá ser,
eventualmente, uma forma indireta muito eficaz de verificar resultados em
Matemática.

A teoria matemática é construída pelo \emph{Modelo Axiomático}, isso
significa que partimos de alguns conceitos primitivos e afirmações básicas
(\emph{axiomas ou postulados}) que aceitamos como verdadeiras, isto é, sem
que sejam demonstradas. Os conceitos de ponto, reta e plano são conceitos
primitivos. Quando afirmamos que ``Por dois pontos passam uma única reta'',
temos um exemplo de um axioma. A partir desses conceitos, todos os outros
conceitos podem ser definidos e todas as outras proposições demonstradas.
Assim, \emph{definições}  conceitos dados em função de termos considerados
previamente conhecidos.

Existem apenas duas categorias de enunciados: as proposições primitivas
(axiomas), que são proposições aceitas sem demonstração (não havendo
preocupação se são evidentes ou não); e as proposições demonstradas (teoremas,
proposições, corolários e lemas) por meio de raciocínios logicamente corretos,
a partir dos postulados e definições pertinentes.

Nesse sentido, entendemos por \emph{demonstração} Matemática, o processo de
raciocínio lógico e dedutivo para checar a veracidade de uma proposição
condicional. Nesse processo são usados argumentos válidos e definições, ou
seja, aqueles que concluam afirmações verdadeiras a partir de fatos que também
são verdadeiros.

As proposições que podem ser demonstradas são chamadas de \emph{Teoremas}.
Em geral, denominamos de Teorema apenas certas proposições que são de grande
importância matemática, chamando-se simplesmente de proposição ao resto das
proposições verdadeiras que admitem uma demonstração. Para que um Teorema seja
verificado, as hipóteses tem que ser cuidadosamente satisfeitas. Um famoso
Teorema é o \emph{Teorema de Pitágoras}, que enunciaremos a seguir:

%------------------------------------------------------------------------------%
\begin{teorema}{Teorema de Pitágoras}{teo:pitagoras}
  Em um triângulo retângulo a
  soma dos quadrados dos catetos é igual ao quadrado da hipotenusa.
\end{teorema}
%------------------------------------------------------------------------------%

Note que o teorema de Pitágoras não está na forma ``Se P então Q''. No entanto,
podemos reescrevê-lo nesse formato.

%------------------------------------------------------------------------------%
\begin{teorema}{Teorema de Pitágoras no Gradado}{teo:pitagoras-quadrado}
  Se $T$ é um triângulo retângulo
  de catetos $b$ e $c$ e hipotenusa $a$, então $a^2=b^2+c^2$.
\end{teorema}
%------------------------------------------------------------------------------%

No Teorema de Pitágoras temos que as hipóteses são ``T é um triângulo retângulo
de catetos $b$ e $c$ e hipotenusa $a$'' e a tese é ``$a^2=b^2+c^2$ ''. Existem
algumas sentenças que são diretas consequências de um determinado teorema são
chamadas de \emph{Corolários}, já resultados auxiliares que são usados para
demonstrar uma determinada proposição ou teorema são chamados de
\emph{Lemas}.

Quando em uma sentença temos que $P \Rightarrow Q$ e $Q \Rightarrow P$,
escrevemos $P \Leftrightarrow Q$ e o símbolo $\Leftrightarrow$ é lido como ``se
e somente se''. Em algumas proposições e teoremas matemáticos vale a recíproca,
mas nem sempre isso é verdade.

%------------------------------------------------------------------------------%
